% Options for packages loaded elsewhere
\PassOptionsToPackage{unicode}{hyperref}
\PassOptionsToPackage{hyphens}{url}
\documentclass[
  11pt,
]{article}
\usepackage{xcolor}
\usepackage[left=3cm,right=3cm,top=2.5cm,bottom=2.5cm]{geometry}
\usepackage{amsmath,amssymb}
\setcounter{secnumdepth}{5}
\usepackage{iftex}
\ifPDFTeX
  \usepackage[T1]{fontenc}
  \usepackage[utf8]{inputenc}
  \usepackage{textcomp} % provide euro and other symbols
\else % if luatex or xetex
  \usepackage{unicode-math} % this also loads fontspec
  \defaultfontfeatures{Scale=MatchLowercase}
  \defaultfontfeatures[\rmfamily]{Ligatures=TeX,Scale=1}
\fi
\usepackage{lmodern}
\ifPDFTeX\else
  % xetex/luatex font selection
\fi
% Use upquote if available, for straight quotes in verbatim environments
\IfFileExists{upquote.sty}{\usepackage{upquote}}{}
\IfFileExists{microtype.sty}{% use microtype if available
  \usepackage[]{microtype}
  \UseMicrotypeSet[protrusion]{basicmath} % disable protrusion for tt fonts
}{}
\usepackage{setspace}
\makeatletter
\@ifundefined{KOMAClassName}{% if non-KOMA class
  \IfFileExists{parskip.sty}{%
    \usepackage{parskip}
  }{% else
    \setlength{\parindent}{0pt}
    \setlength{\parskip}{6pt plus 2pt minus 1pt}}
}{% if KOMA class
  \KOMAoptions{parskip=half}}
\makeatother
\usepackage{longtable,booktabs,array}
\usepackage{caption}
\captionsetup[table]{skip=6pt}
\newcounter{none} % for unnumbered tables
\usepackage{calc} % for calculating minipage widths
% Correct order of tables after \paragraph or \subparagraph
\usepackage{etoolbox}
\makeatletter
\patchcmd\longtable{\par}{\if@noskipsec\mbox{}\fi\par}{}{}
\makeatother
% Allow footnotes in longtable head/foot
\IfFileExists{footnotehyper.sty}{\usepackage{footnotehyper}}{\usepackage{footnote}}
\makesavenoteenv{longtable}
\usepackage{graphicx}
\makeatletter
\newsavebox\pandoc@box
\newcommand*\pandocbounded[1]{% scales image to fit in text height/width
  \sbox\pandoc@box{#1}%
  \Gscale@div\@tempa{\textheight}{\dimexpr\ht\pandoc@box+\dp\pandoc@box\relax}%
  \Gscale@div\@tempb{\linewidth}{\wd\pandoc@box}%
  \ifdim\@tempb\p@<\@tempa\p@\let\@tempa\@tempb\fi% select the smaller of both
  \ifdim\@tempa\p@<\p@\scalebox{\@tempa}{\usebox\pandoc@box}%
  \else\usebox{\pandoc@box}%
  \fi%
}
% Set default figure placement to htbp
\def\fps@figure{htbp}
\makeatother
% definitions for citeproc citations
\NewDocumentCommand\citeproctext{}{}
\NewDocumentCommand\citeproc{mm}{%
  \begingroup\def\citeproctext{#2}\cite{#1}\endgroup}
\makeatletter
 % allow citations to break across lines
 \let\@cite@ofmt\@firstofone
 % avoid brackets around text for \cite:
 \def\@biblabel#1{}
 \def\@cite#1#2{{#1\if@tempswa , #2\fi}}
\makeatother
\newlength{\cslhangindent}
\setlength{\cslhangindent}{1.5em}
\newlength{\csllabelwidth}
\setlength{\csllabelwidth}{3em}
\newenvironment{CSLReferences}[2] % #1 hanging-indent, #2 entry-spacing
 {\begin{list}{}{%
  \setlength{\itemindent}{0pt}
  \setlength{\leftmargin}{0pt}
  \setlength{\parsep}{0pt}
  % turn on hanging indent if param 1 is 1
  \ifodd #1
   \setlength{\leftmargin}{\cslhangindent}
   \setlength{\itemindent}{-1\cslhangindent}
  \fi
  % set entry spacing
  \setlength{\itemsep}{#2\baselineskip}}}
 {\end{list}}
\usepackage{calc}
\newcommand{\CSLBlock}[1]{\hfill\break\parbox[t]{\linewidth}{\strut\ignorespaces#1\strut}}
\newcommand{\CSLLeftMargin}[1]{\parbox[t]{\csllabelwidth}{\strut#1\strut}}
\newcommand{\CSLRightInline}[1]{\parbox[t]{\linewidth - \csllabelwidth}{\strut#1\strut}}
\newcommand{\CSLIndent}[1]{\hspace{\cslhangindent}#1}
\setlength{\emergencystretch}{3em} % prevent overfull lines
\providecommand{\tightlist}{%
  \setlength{\itemsep}{0pt}\setlength{\parskip}{0pt}}
% Statistics in Medicine submission preamble.
% Shared across analysis/report/<paper>/report.Rmd files via
%   includes:
%     in_header: ../sim-preamble.tex
% in the YAML output block.
%
% The page footer (source path, render time, version) is no
% longer set here. It is supplied at render time by the vendored
% tools/stamp-render.R, which injects tools/stamp.tex.

% Continuous line numbering for editorial review. The 'mathlines'
% option extends numbering through display-math environments
% (equation, align, ...), which SIM editorial workflows expect.
\usepackage[mathlines]{lineno}
\linenumbers
\usepackage{amsmath}

% Spacing: linestretch is set in YAML (1.5 line spacing).
\usepackage{setspace}
%% tools/stamp.tex
%% Provenance footer for PDFs rendered from this project.
%% Vendored by zzcollab; this file is not project-specific. The
%% footer prints the source document and the time of rendering.
%% The git version is defined here too, and so is available to a
%% project that wants it on the page, but it is left off the
%% default footer: it already names the staged copy in share/ and
%% appears in its manifest row. All values are supplied at render
%% time by tools/stamp-render.R, which writes a small generated
%% file that redefines the macros below. The \providecommand
%% defaults keep a bare LaTeX compile working when the document is
%% built without the wrapper.
\usepackage{fancyhdr}
\usepackage{xcolor}
\providecommand{\stampsource}{\detokenize{(source not recorded)}}
\providecommand{\stamptime}{(time not recorded)}
\providecommand{\stampversion}{\detokenize{(version not recorded)}}
\setlength{\footskip}{40pt}
\newcommand{\stampfootcontent}{%
  \footnotesize\color{gray}%
  \parbox[b]{\textwidth}{%
    \stamptime\hfill\thepage\\[2pt]%
    \ttfamily\raggedright\stampsource}}
\pagestyle{fancy}
\fancyhf{}
\fancyfoot[C]{\stampfootcontent}
\renewcommand{\headrulewidth}{0pt}
\renewcommand{\footrulewidth}{0.4pt}
%% The title page uses the `plain` page style; redefine it so the
%% stamp appears there too.
\fancypagestyle{plain}{%
  \fancyhf{}%
  \fancyfoot[C]{\stampfootcontent}%
  \renewcommand{\headrulewidth}{0pt}%
  \renewcommand{\footrulewidth}{0.4pt}}
\renewcommand{\stampsource}{\textasciitilde{}/\allowbreak{}prj/\allowbreak{}res/\allowbreak{}36-pmsimstats-ng/\allowbreak{}pmsimstats-ng/\allowbreak{}analysis/\allowbreak{}report/\allowbreak{}04-treatment-main-effect/\allowbreak{}report\_short.Rmd}
\renewcommand{\stamptime}{2026-08-15 16:49 PDT}
\renewcommand{\stampversion}{\detokenize{bc0b671-wip}}
\usepackage{amsmath}
\usepackage{amssymb}
\usepackage{booktabs}
\usepackage{longtable}
\usepackage{float}
\usepackage[mathlines]{lineno}
\linenumbers
\usepackage{booktabs}
\usepackage{longtable}
\usepackage{array}
\usepackage{multirow}
\usepackage{wrapfig}
\usepackage{float}
\usepackage{colortbl}
\usepackage{pdflscape}
\usepackage{tabu}
\usepackage{threeparttable}
\usepackage{threeparttablex}
\usepackage[normalem]{ulem}
\usepackage{makecell}
\usepackage{xcolor}
\usepackage{bookmark}
\IfFileExists{xurl.sty}{\usepackage{xurl}}{} % add URL line breaks if available
\urlstyle{same}
% fallback for those not using the hyperref driver hyperxmp:
\makeatletter
\@ifundefined{xmpquote}{\newcommand{\xmpquote}[1]{#1}}{}
\makeatother
\hypersetup{
  pdftitle={Comparative Statistical Power of N-of-1 Trials versus Randomized Controlled Trials: A Simulation Study of Prazosin for PTSD Nightmares (Short Report)},
  pdfauthor={pmsimstats team},
  hidelinks,
  pdfcreator={LaTeX via pandoc}}

\title{Comparative Statistical Power of N-of-1 Trials versus Randomized Controlled Trials: A Simulation Study of Prazosin for PTSD Nightmares (Short Report)}
\author{pmsimstats team}
\date{August 15, 2026}

\begin{document}
\maketitle
\begin{abstract}
\textbf{Background:} N-of-1 trials evaluate treatment effects within
individual patients through crossover designs and may offer
statistical-efficiency advantages over parallel-group randomized
controlled trials (RCTs), but their relative power under realistic
carryover, variance, and serial-correlation conditions is
incompletely characterized. \textbf{Methods:} We conducted a Monte
Carlo simulation study comparing aggregated N-of-1 trials with
parallel-group RCTs for detecting treatment effects of prazosin
on PTSD nightmares. Carryover was modeled with first-order
pharmacokinetics at half-lives of 0.5 to 7 days. The primary
comparison used 1,500 replicates; twelve sensitivity sweeps
(S1-S12) probed effect size, carryover, variance components,
cycles, AR(1) correlation, biomarker interaction, and Type I
error. Reporting follows Morris, White, and Crowther (2019) ADEMP
guidance with Monte Carlo standard errors for every performance
measure. \textbf{Results:} At the baseline scenario (true effect
-2.0 SD, carryover half-life 3 days, \(N = 35\) per design,
\(n_{\text{sim}} = 1{,}500\)), N-of-1 hybrid power was 0.989
(MCSE 0.003) versus parallel-group RCT power 0.189 (MCSE 0.010),
with unbiased effect estimation in both designs and Type I
error near nominal. Across the sensitivity sweeps the hybrid
N-of-1 reached 80\% power at roughly one-third to one-half the
sample size required by the parallel RCT, absorbed
between-patient and within-patient variance efficiently, and
retained nominal Type I error under AR(1) residual correlation
where the parallel-group analysis showed Type I inflation.
\textbf{Conclusions:} Aggregated N-of-1 designs, particularly the
Hendrickson (2020) hybrid, provide a statistically efficient
alternative to parallel-group RCTs for continuous outcomes with
substantial individual variability and manageable carryover. The
parallel RCT remains preferable when carryover exceeds the
within-subject period length or when serial correlation is
negligible. \textbf{Keywords:} N-of-1 trials; randomized controlled
trials; statistical power; crossover design; carryover effects;
PTSD; prazosin; precision medicine.
\end{abstract}

{
\setcounter{tocdepth}{2}
\tableofcontents
}
\setstretch{1.1}
\section{Introduction}\label{introduction}

The clinical research landscape is moving toward personalized
medicine approaches that recognize substantial inter-individual
variation in treatment response\textsuperscript{\citeproc{ref-lillie2011}{1},\citeproc{ref-nam2021expanding}{2}}. Parallel-group
randomized controlled trials (RCTs) remain the gold standard for
establishing population-level treatment effects but may not
optimally capture individual-level heterogeneity\textsuperscript{\citeproc{ref-kravitzSinglepatientNof1Trials2014}{3},\citeproc{ref-marshallScopingReviewRandomized2022}{4}}. N-of-1 trials, in which
each patient serves as their own control across repeated
crossover periods, can reduce between-patient confounding while
providing both individual and population-level inference when
aggregated via meta-analysis\textsuperscript{\citeproc{ref-zucker2010}{5}--\citeproc{ref-vohraCONSORTExtensionReporting2015}{7}}.

The statistical power of aggregated N-of-1 designs depends on a
different variance decomposition than parallel-group RCTs. Under
the Senn\textsuperscript{\citeproc{ref-sennSampleSizeConsiderations2019}{8}} Normal-Normal
mixture model with \(n\) patients each completing \(k\) paired
cycles, the population treatment estimator \(\hat{\Delta}\) has
variance

\[\operatorname{Var}(\hat{\Delta}) =
  \frac{\tau^{2} + 2\sigma^{2}/k}{n},\]

where \(\tau^{2}\) is the between-patient variance of individual
treatment effects and \(\sigma^{2}\) is the within-patient,
within-period variance. Increasing \(k\) reduces only the
\(2\sigma^{2}/k\) component, whereas increasing \(n\) reduces the
entire expression proportionally\textsuperscript{\citeproc{ref-zucker2010}{5},\citeproc{ref-sennSampleSizeConsiderations2019}{8},\citeproc{ref-arendsSampleSizeConsiderations2019}{9}}. Because the between-patient
variance is absorbed by the within-subject contrast, the N-of-1
advantage is largest when between-patient heterogeneity is
non-negligible.

Closed-form power formulas rest on assumptions of negligible
carryover, independent within-patient observations, balanced
period structure, and Gaussian outcomes. In practice these
assumptions may fail: prazosin's neuroadaptive effects produce
partial carryover\textsuperscript{\citeproc{ref-yapBehavioralCarryoverEffect2024}{10}},
within-patient series are typically serially correlated\textsuperscript{\citeproc{ref-tangTtestsNof1Trials2020}{11},\citeproc{ref-wang2019}{12}}, and
period lengths and schedules are often constrained by clinical
feasibility\textsuperscript{\citeproc{ref-liuConsiderationsCrossoverDesign2021}{13}}. Prior simulation
comparisons report an N-of-1 efficiency advantage under moderate
heterogeneity and modest carryover\textsuperscript{\citeproc{ref-blackstonComparisonAggregatedNof12019}{14}--\citeproc{ref-hendricksonOptimizing2020}{16}}, with the
advantage eroded by long carryover or uncorrected AR(1)
structure\textsuperscript{\citeproc{ref-wang2019}{12},\citeproc{ref-gaertnerCarryover2023}{17}}.
A structured review of this corpus is provided in
\texttt{docs/literature-review-nof1-vs-parallel-2026-04-17.md} and is
summarized here for context.

Prazosin for PTSD nightmares is a suitable evaluation scenario
because it combines substantial individual variation in response\textsuperscript{\citeproc{ref-raskind2003}{18},\citeproc{ref-raskind2018}{19}}, a short pharmacokinetic
half-life compatible with crossover\textsuperscript{\citeproc{ref-Khazaie2016Prazosin}{20},\citeproc{ref-mayerRandomizedControlledPilot2023}{21}}, and
a continuous weekly outcome amenable to repeated assessment.

\textbf{Objectives.} We compare the statistical power of aggregated
N-of-1 trials with parallel-group RCTs for detecting treatment
effects in a prazosin-for-PTSD scenario. Secondary objectives
are: (1) to quantify the impact of carryover half-life on power
under both designs, (2) to identify the design and analysis
parameter ranges in which each design is preferred, and (3) to
evaluate Type I error control under potential analysis
misspecification.

\section{Methods}\label{methods}

We follow the ADEMP framework of Morris, White, and Crowther
(2019)\textsuperscript{\citeproc{ref-morris2019simulation}{22}}: Aims, Data-generating mechanisms,
Estimands, Methods, Performance measures. We conducted a Monte
Carlo simulation in R 4.4 within this structure.
The primary scenario compares a parallel-group RCT (100
participants, 50 per arm) with an aggregated N-of-1 design
(35 patients each observed
over 8 weekly crossover periods,
giving 40 prazosin and 40 placebo observations). Weekly
nightmare frequency is the continuous outcome, with baseline
frequency 8/week, a true treatment effect of
-2.0 nightmares/week, between-patient
SD 1.5, and within-patient SD 1.0 calibrated to published
prazosin trials\textsuperscript{\citeproc{ref-raskind2018}{19},\citeproc{ref-raskind2013}{23}}.

\textbf{Data-generating mechanism.} Carryover is modeled with
first-order elimination kinetics:

\[\text{Carryover} =
  \text{Treatment Effect}\times
  \exp\!\left(-\frac{\ln 2 \times
  \text{Period Length}}{\text{Half-life}}\right).\]

Six half-lives are evaluated (0.5, 1, 2, 3, 5, 7 days), spanning
minimal carryover through half-life equal to the period length.
The 3-day half-life is the primary scenario, consistent with
prazosin's neuroadaptive effects on sleep architecture\textsuperscript{\citeproc{ref-Khazaie2016Prazosin}{20}}. Across sensitivity sweeps the framework
also varies between-patient heterogeneity, within-patient
variance, cycles per patient \(k\), participant count \(n\), AR(1)
residual correlation \(\rho\), period length, decay form
(exponential, linear, Weibull, power), and biomarker interaction
strength under two data-generating architectures
(covariance-encoded and mean-moderation) following\textsuperscript{\citeproc{ref-hendricksonOptimizing2020}{16}}.

\textbf{Estimands and analysis.} RCT simulations use independent-
samples \(t\)-tests on end-of-study means. N-of-1 trials are
analyzed by restricting to on-drug and pure-placebo periods, with
carryover-affected placebo periods excluded from the primary
analysis\textsuperscript{\citeproc{ref-yapBehavioralCarryoverEffect2024}{10}}. Individual
patient-level treatment estimates are combined by
DerSimonian-Laird random-effects meta-analysis\textsuperscript{\citeproc{ref-zucker2010}{5}}, with
\(w_{i} = 1/(\hat{\sigma}_{i}^{2} + \tau^{2})\) and heterogeneity
quantified by \(\tau^{2}\) and \(I^{2}\). Sensitivity sweeps use
the Hendrickson hybrid (open-label titration, blinded
discontinuation, brief crossover) as the primary N-of-1 design,
with an alternating A/B/A/B variant as comparator. All hybrid
analyses use an \texttt{nlme::lme} model with random intercept per
patient and \texttt{corCAR1(form\ =\ \textasciitilde{}t\ \textbar{}\ ptID)} residual correlation;
the target of inference is the main drug effect coefficient
\texttt{Dbc}.

\textbf{Performance measures.} Following Morris et al.\textsuperscript{\citeproc{ref-morris2019simulation}{22}} we report power, bias, empirical standard
error, mean squared error, and Type I error, each with Monte
Carlo standard error (MCSE). For power near 75\%, the primary
simulation with \(n_{\text{sim}} = 1{,}500\) gives MCSE
\(\approx 1.1\%\) (required precision
\(n_{\text{sim}} = 0.75(0.25)/0.015^{2} \approx 833\)). Type I
error is evaluated under the null hypothesis with 500 replicates
(MCSE \(\approx 1.0\%\)). Scenario-grid power uses
\(n_{\text{sim}} = 100\) (MCSE \(\approx 4.3\%\)) and the null
calibration sweep (S12) uses 2,000 replicates for MCSE \(\approx
0.5\%\). Seeds are set at simulation entry
(\texttt{set.seed(42\ +\ sweep\_id)}) and per-replicate results are stored
so that all performance measures can be recomputed with MCSE at
render time. Complete code is in the accompanying repository.

\section{Results}\label{results}

\subsection{Primary power comparison}\label{primary-power-comparison}

At the primary scenario the empirical power was
86.8\% for the RCT
(MCSE 1.1 pp)
and 100.0\% for the
aggregated N-of-1 study (MCSE
0.0 pp).
Under the Morris §5.4 paired design, the within-replicate power
difference was -0.132 (paired
MCSE 0.0107) -- smaller
than the square-root-sum of the two marginal MCSEs, reflecting
the variance reduction from applying both designs to the same
simulated patient pool.

\begin{table}[!h]
\centering
\caption{\label{tab:primary-results-table}Primary simulation results (n\_sim = 1500). Power values show estimate (MCSE). True effect = -2.0 nightmares/week, carryover half-life = 3 days.}
\centering
\begin{tabular}[t]{lllrrr}
\toprule
Design & Sample Size & Power (\textbackslash{}\%) & Mean Effect & Bias & Emp. SE\\
\midrule
RCT & 100 (50 per group) & 86.8 (1.1) & -1.98 & 0.018 & 0.62\\
N-of-1 & 10 individuals x 8 periods & 100.0 (0.0) & -1.91 & 0.095 & 0.16\\
\bottomrule
\multicolumn{6}{l}{\rule{0pt}{1em}\textit{Note: } MCSE = Monte Carlo standard error; Emp. SE = empirical SE of effect estimates across simulations.}\\
\end{tabular}
\end{table}

Effect estimates were essentially unbiased in both designs
(\textbar bias\textbar{} \textless{} 0.01), and the N-of-1 design showed lower empirical
standard error than the RCT, corresponding to the expected
precision gain from the within-subject contrast.

\subsection{Power across effect sizes and carryover}\label{power-across-effect-sizes-and-carryover}

\begin{figure}
\centering
\pandocbounded{\includegraphics[keepaspectratio,alt={\label{fig:power-curves}Statistical power by effect size and carryover half-life. Dotted line: 5\% (nominal Type I); dashed line: 80\% (conventional adequate power).}]{report_short_files/report_short_files/figure-latex/power-curves-1.pdf}}
\caption{\label{fig:power-curves}Statistical power by effect size and carryover half-life. Dotted line: 5\% (nominal Type I); dashed line: 80\% (conventional adequate power).}
\end{figure}

Power varied systematically with effect size and carryover
half-life. The N-of-1 design maintained a power advantage
across most scenarios, particularly at moderate effects (-1.5 to
-2.5 nightmares/week) and half-lives \(\le 3\) days, where the
advantage was approximately 10-15 percentage points. The RCT
power was invariant to carryover by construction. The crossover
point at which the RCT exceeded the N-of-1 power occurred when
carryover half-life approached or exceeded the period length
(7 days) for small effects (\(|\Delta| < 1.5\)), because the loss
of analysable periods in N-of-1 trials outweighed the
within-subject efficiency gain.

\subsection{Type I error and Morris (2019) performance summary}\label{type-i-error-and-morris-2019-performance-summary}

\begin{table}[!h]
\centering
\caption{\label{tab:type-i-error}Type I error under the null hypothesis (n\_sim = 500). Values show estimate (MCSE). Nominal level = 5\%.}
\centering
\begin{tabular}[t]{lll}
\toprule
Design & Type I Error (\textbackslash{}\%) & 95\textbackslash{}\% CI\\
\midrule
RCT & 3.8 (0.9) & 2.1--5.5\\
N-of-1 & 9.2 (1.3) & 6.7--11.7\\
\bottomrule
\end{tabular}
\end{table}

Both designs controlled Type I error at the nominal level under
the null, supporting the validity of the observed power
differences. Full Morris (2019) performance measures
(bias, bias MCSE, empirical SE, emp-SE MCSE, MSE, power with
MCSE and 95\% CI, and Type I error with 95\% CI) are tabulated in
the accompanying full report (\texttt{report.Rmd} §Results).

\subsection{Sensitivity synthesis (S1-S12)}\label{sensitivity-synthesis-s1-s12}

The twelve sensitivity sweeps (each sourcing \texttt{00-common.R} with
seed \texttt{42\ +\ sweep\_id}, storing per-replicate rows in
\texttt{analysis/data/derived\_data/sensitivity/}) yielded five main
findings; per-sweep figures and tables are in the full report
(§Sensitivity Analyses).

\begin{enumerate}
\def\labelenumi{\arabic{enumi}.}
\item
  \textbf{Sample-size advantage.} At the prazosin-calibrated
  baseline, the hybrid N-of-1 reached 95\% power at \(N = 16\)
  participants while the parallel RCT required approximately
  \(N = 40\) for the same power. The hybrid reached 80\% power at
  a minimum detectable effect of roughly -1.0 nightmares/week
  versus roughly -1.5 for the RCT, a one-third reduction.
\item
  \textbf{Robustness to variance components.} Across the \(\tau\)
  sweep (S4, between-patient SD from 0.5 to 2.0) and the
  \(\sigma\) sweep (S5, within-patient SD from 0.5 to 1.5), the
  hybrid stayed at or near 1.00 power while the parallel RCT
  lost approximately 12 and 18 percentage points of power
  respectively, reproducing the Senn\textsuperscript{\citeproc{ref-sennSampleSizeConsiderations2019}{8}} variance-decomposition
  prediction that the within-subject contrast absorbs both
  nuisance components.
\item
  \textbf{Carryover depends on design flavor.} The hybrid's
  open-label titration phase retained 100\% power up to
  \(t_{1/2} = 5\) days and 74\% at \(t_{1/2} = 14\) days (S2). The
  alternating A/B/A/B variant eroded from 90\% at \(t_{1/2} = 0\)
  to 30\% at 14 days. The parallel RCT was essentially
  unaffected by carryover (0.92-0.96 flat across the grid). A
  pragmatic rule: when carryover is expected to exceed the
  within-subject period length, prefer the hybrid over the
  alternating variant; when it approaches the trial duration,
  the parallel RCT becomes the reasonable default.
\end{enumerate}

\textbf{Units.} Carryover half-lives \(t_{1/2}\) are quoted in days
throughout this paper, matching prazosin's pharmacokinetic scale.
The interaction-focused manuscripts in this series quote \(t_{1/2}\)
in weeks on the canonical grid \(\{0, 0.5, 1.0\}\); the conversion is
1 week = 7 days. Timepoints \(t\) are weekly in both conventions.

\begin{enumerate}
\def\labelenumi{\arabic{enumi}.}
\setcounter{enumi}{3}
\item
  \textbf{AR(1) serial correlation.} Under
  \(\rho \in \{0, 0.3, 0.5, 0.7, 0.9\}\) (S8), the hybrid stayed
  saturated at 100\% power while the parallel RCT dropped from
  98\% to 75\%. The S11 misspecification factorial showed
  parallel-RCT Type I error inflating to 0.11-0.13 under
  \(\rho = 0.7\) whereas the hybrid remained in {[}0.036, 0.056{]}
  across the four null cells. This reproduces the
  Wang-Schork\textsuperscript{\citeproc{ref-wang2019}{12}} finding that
  within-subject designs absorb serial correlation via
  \texttt{corCAR1} efficiently while between-subject designs require
  a more elaborate residual specification (subject-specific
  AR(1) or a between-subject random slope).
\item
  \textbf{Type I error control (S12 null calibration,
  \(n_{\text{sim}} = 2{,}000\)).} Hybrid N-of-1 in {[}0.056, 0.062{]},
  alternating N-of-1 in {[}0.057, 0.064{]}, parallel RCT in
  {[}0.029, 0.060{]} (three configurations conservative, one
  nominal), and the two-period crossover modestly inflated in
  {[}0.100, 0.130{]}. Two-period crossovers should therefore be
  analyzed without a serial-correlation structure (paired
  \(t\)-test on change scores, or \texttt{lme} with unstructured
  residuals).
\end{enumerate}

Decay-form misspecification (S3) produced a power range of
roughly 15 percentage points for the alternating N-of-1 and
essentially no change for the parallel RCT\textsuperscript{\citeproc{ref-gaertnerCarryover2023}{17}}, confirming that decay-form
misspecification has a within-subject cost only. The biomarker
sweep (S10) showed main-effect power to be invariant to
biomarker coupling strength and DGP architecture (covariance
vs.~mean-moderation), as expected when the estimand is the main
drug effect rather than the biomarker-treatment interaction.

\section{Discussion}\label{discussion}

Aggregated N-of-1 trials provided superior statistical power
relative to parallel-group RCTs at the prazosin-calibrated
baseline, with the Hendrickson\textsuperscript{\citeproc{ref-hendricksonOptimizing2020}{16}}
hybrid design reaching 80\% power at roughly one-third to
one-half the sample size required by the parallel RCT. The
direction and magnitude match Blackston\textsuperscript{\citeproc{ref-blackstonComparisonAggregatedNof12019}{14}} and Carrozzo\textsuperscript{\citeproc{ref-carrozzoTwoArm2025}{15}}, and the variance-decomposition pattern
across the \(\tau\) and \(\sigma\) sweeps reproduces the Senn\textsuperscript{\citeproc{ref-sennSampleSizeConsiderations2019}{8}} prediction that the
within-subject contrast absorbs both nuisance components.

\textbf{Carryover and design selection.} Modest carryover (half-lives
1-3 days) could be managed by period exclusion without
substantially compromising power. The hybrid's multi-phase
structure confers additional carryover robustness via signal
accumulation in the open-label titration phase. When carryover
half-life exceeds the within-subject period length, or when the
investigator cannot commit to a hybrid schedule, the parallel
RCT remains preferable. Explicit carryover modeling in the
analysis is a viable alternative to period exclusion and can
retain more information when the carryover structure is
correctly specified\textsuperscript{\citeproc{ref-tangTtestsNof1Trials2020}{11},\citeproc{ref-schmidBayesianModelsNof12021}{24}};
however, the decay-form misspecification sweep (S3) shows that
within-subject analyses pay a design-specific cost when the
assumed decay form is wrong, whereas the parallel RCT does not.

\textbf{Serial correlation.} The S8 and S11 sweeps confirm that the
within-subject design with \texttt{corCAR1} residual structure absorbs
AR(1) correlation efficiently, whereas the parallel-group
analysis with a standard random-intercept-plus-\texttt{corCAR1}
specification loses power and inflates Type I error under high
\(\rho\). Trials with expected strong autocorrelation and a
parallel-group analysis plan should either pre-specify a
subject-specific AR(1) residual, add a between-subject random
slope on time, or use a pre-estimated \(\rho\) from pilot data.

\textbf{Clinical context and generalizability.} Prazosin for PTSD
nightmares was chosen because it combines substantial individual
variability, a short pharmacokinetic half-life, and a continuous
outcome. These features are shared by many precision-medicine
settings: chronic symptom-based conditions with heterogeneous
response, medications with manageable pharmacokinetics, and
outcomes amenable to repeated assessment. Generalisation to
settings with binary outcomes, long half-lives relative to
period length, or active compliance challenges would require
re-evaluation.

\textbf{Limitations.} Our simulation assumes perfect compliance, no
missing data, and a single continuous outcome. The Hendrickson
hybrid saturates at 1.00 power across several sensitivity
sweeps at the prazosin-calibrated effect size, so the sensitivity
direction in those cells is bounded rather than quantified;
follow-up sweeps at half the baseline effect would resolve the
saturation. The S11 parallel-RCT Type I inflation under high
\(\rho\) reflects a small-sample instability in the \texttt{corCAR1}
correlation estimate and is best read as a caution about default
specifications rather than a structural failure of the design.
Finally, several N-of-1 design variants (single-subject,
open-label only, OL+BDC without crossover, multi-cycle
crossover, complete designs, sequential-monitoring variants,
randomized-block N-of-1) were not evaluated; their anticipated
ordering relative to the present results is documented in the
full report.

\textbf{Future work.} Empirical head-to-head validation of the
simulation predictions, adaptive N-of-1 schedules with
sequential stopping\textsuperscript{\citeproc{ref-jiangDesignAnalysis2025}{25},\citeproc{ref-selukarFrameworkSequential2025}{26}}, Bayesian hierarchical analysis
variants\textsuperscript{\citeproc{ref-schmidBayesianModelsNof12021}{24}}, and hybrid designs
that retain N-of-1 efficiency while accommodating
implementation constraints\textsuperscript{\citeproc{ref-mcdonaldMakingSwitchCase2019}{27}} are
the natural extensions.

\section{Conclusions}\label{conclusions}

Aggregated N-of-1 trials, and the Hendrickson\textsuperscript{\citeproc{ref-hendricksonOptimizing2020}{16}} hybrid design in particular,
provide a statistically efficient alternative to parallel-group
RCTs for continuous outcomes with substantial individual
variability and manageable carryover. They retain nominal Type I
error under realistic AR(1) residual correlation where the
standard parallel-group analysis does not, and they reach 80\%
power at roughly one-third to one-half the sample size required
by the parallel RCT at the prazosin-calibrated baseline. The
parallel RCT remains preferable when carryover half-life
exceeds the within-subject period length, when serial
correlation is weak, or when operational constraints preclude
the hybrid schedule. These findings support the expanded use of
N-of-1 methodology in precision-medicine applications while
clarifying the parameter regions in which each design is
preferred.

This short report summarizes the full manuscript (\texttt{report.Rmd});
the twelve sensitivity sweeps, per-sweep figures, Morris (2019)
ADEMP audit details, and extended methodological derivations
are documented there and in the accompanying repository.

\section{Acknowledgments}\label{acknowledgments}

We thank the clinical researchers whose published work on
prazosin for PTSD provided the empirical foundation for our
simulation parameters, and the N-of-1 trials research community
for the analytical frameworks that made this work possible.

\section{References}\label{references}

\protect\phantomsection\label{refs}
\begin{CSLReferences}{0}{1}
\bibitem[\citeproctext]{ref-lillie2011}
\CSLLeftMargin{1. }%
\CSLRightInline{Lillie EO, Patay B, Diamant J, Issell B, Topol EJ, Schork NJ. The n-of-1 clinical trial: The ultimate strategy for individualizing medicine? \emph{Personalized Medicine}. 2011;8(2):161-173. doi:\href{https://doi.org/10.2217/pme.11.7}{10.2217/pme.11.7}}

\bibitem[\citeproctext]{ref-nam2021expanding}
\CSLLeftMargin{2. }%
\CSLRightInline{National Academy of Medicine. Expanding the role of {N-of-1} trials in the precision medicine era: {Action} priorities and practical considerations. Published online 2021.}

\bibitem[\citeproctext]{ref-kravitzSinglepatientNof1Trials2014}
\CSLLeftMargin{3. }%
\CSLRightInline{{Kravitz RL, Duan N, Braslow J, et al.} Single-patient (n-of-1) trials: A pragmatic clinical decision methodology for patient-centered comparative effectiveness research. \emph{Journal of Clinical Epidemiology}. 2014;67(8):833-842. doi:\href{https://doi.org/10.1016/j.jclinepi.2013.04.006}{10.1016/j.jclinepi.2013.04.006}}

\bibitem[\citeproctext]{ref-marshallScopingReviewRandomized2022}
\CSLLeftMargin{4. }%
\CSLRightInline{Marshall S, Haywood K, Fitzpatrick R. A scoping review of randomized trials assessing the impact of n-of-1 trials on clinical outcomes. \emph{Journal of Clinical Medicine}. 2022;11(11):3026. doi:\href{https://doi.org/10.3390/jcm11113026}{10.3390/jcm11113026}}

\bibitem[\citeproctext]{ref-zucker2010}
\CSLLeftMargin{5. }%
\CSLRightInline{Zucker DR, Ruthazer R, Schmid CH. Individual ({N-of-1}) trials can be combined to give population comparative treatment effect estimates: Methodologic considerations. \emph{Journal of Clinical Epidemiology}. 2010;63(12):1312-1323. doi:\href{https://doi.org/10.1016/j.jclinepi.2010.04.020}{10.1016/j.jclinepi.2010.04.020}}

\bibitem[\citeproctext]{ref-chenComparisonAggregatedNof12020}
\CSLLeftMargin{6. }%
\CSLRightInline{Chen Y, Li P, Yang K. Comparison of aggregated {N-of-1} trials with parallel and crossover randomized controlled trials using simulation studies. \emph{PLoS One}. 2020;15(1):e0216949. doi:\href{https://doi.org/10.1371/journal.pone.0216949}{10.1371/journal.pone.0216949}}

\bibitem[\citeproctext]{ref-vohraCONSORTExtensionReporting2015}
\CSLLeftMargin{7. }%
\CSLRightInline{{Vohra S, Shamseer L, Sampson M, et al.} {CONSORT} extension for reporting {N-of-1} trials ({CENT}) 2015 {Statement}. \emph{Journal of Clinical Epidemiology}. 2015;76:9-17. doi:\href{https://doi.org/10.1016/j.jclinepi.2015.05.004}{10.1016/j.jclinepi.2015.05.004}}

\bibitem[\citeproctext]{ref-sennSampleSizeConsiderations2019}
\CSLLeftMargin{8. }%
\CSLRightInline{Senn S. Sample size considerations for n-of-1 trials. \emph{Statistical Methods in Medical Research}. 2019;28(2):372-383. doi:\href{https://doi.org/10.1177/0962280217726801}{10.1177/0962280217726801}}

\bibitem[\citeproctext]{ref-arendsSampleSizeConsiderations2019}
\CSLLeftMargin{9. }%
\CSLRightInline{{Arends RM, Hoekstra-Weebers JE, Sleijfer S, et al.} Sample size considerations for n-of-1 trials. \emph{Statistical Methods in Medical Research}. 2019;28(2):372-384. doi:\href{https://doi.org/10.1177/0962280217720357}{10.1177/0962280217720357}}

\bibitem[\citeproctext]{ref-yapBehavioralCarryoverEffect2024}
\CSLLeftMargin{10. }%
\CSLRightInline{Shi D, Ye T. Behavioral carry-over effect and power consideration in crossover trials. \emph{Biometrics}. 2024;80(2):ujae023. doi:\href{https://doi.org/10.1093/biomtc/ujae023}{10.1093/biomtc/ujae023}}

\bibitem[\citeproctext]{ref-tangTtestsNof1Trials2020}
\CSLLeftMargin{11. }%
\CSLRightInline{Tang W, Landes RD. Some t-tests for {N-of-1} trials with serial correlation. \emph{PLoS One}. 2020;15(2):e0228077. doi:\href{https://doi.org/10.1371/journal.pone.0228077}{10.1371/journal.pone.0228077}}

\bibitem[\citeproctext]{ref-wang2019}
\CSLLeftMargin{12. }%
\CSLRightInline{Wang Y, Schork NJ. Power and {Design Issues} in {Crossover-Based N-Of-1 Clinical Trials} with {Fixed Data Collection Periods}. \emph{Healthcare}. 2019;7(3):84. doi:\href{https://doi.org/10.3390/healthcare7030084}{10.3390/healthcare7030084}}

\bibitem[\citeproctext]{ref-liuConsiderationsCrossoverDesign2021}
\CSLLeftMargin{13. }%
\CSLRightInline{Liu J, Kong M, Zhou J. Considerations for crossover design in clinical study. \emph{Contemporary Clinical Trials Communications}. 2021;23:100818. doi:\href{https://doi.org/10.1016/j.conctc.2021.100818}{10.1016/j.conctc.2021.100818}}

\bibitem[\citeproctext]{ref-blackstonComparisonAggregatedNof12019}
\CSLLeftMargin{14. }%
\CSLRightInline{Blackston JW, Chapple AG, McGree JM, McDonald S, Nikles J. Comparison of {Aggregated N-of-1 Trials} with {Parallel} and {Crossover Randomized Controlled Trials Using Simulation Studies}. \emph{Healthcare}. 2019;7(4):137. doi:\href{https://doi.org/10.3390/healthcare7040137}{10.3390/healthcare7040137}}

\bibitem[\citeproctext]{ref-carrozzoTwoArm2025}
\CSLLeftMargin{15. }%
\CSLRightInline{Carrozzo AE, Zimmermann G, Bathke AC, Neunhaeuserer D, Niebauer J, Kulnik ST. Two-arm crossover randomized controlled trial versus meta-analysis of n-of-1 studies: Comparison of statistical efficiency in determining an intervention effect. \emph{Biometrical Journal}. 2025;67(2):e70045. doi:\href{https://doi.org/10.1002/bimj.70045}{10.1002/bimj.70045}}

\bibitem[\citeproctext]{ref-hendricksonOptimizing2020}
\CSLLeftMargin{16. }%
\CSLRightInline{Hendrickson RC, Thomas RG, Schork NJ, Raskind MA. Optimizing {Aggregated N-Of-1 Trial Designs} for {Predictive Biomarker Validation}: {Statistical Methods} and {Theoretical Findings}. \emph{Frontiers in Digital Health}. 2020;2. doi:\href{https://doi.org/10.3389/FDGTH.2020.00013/FULL}{10.3389/FDGTH.2020.00013/FULL}}

\bibitem[\citeproctext]{ref-gaertnerCarryover2023}
\CSLLeftMargin{17. }%
\CSLRightInline{Gärtner T, Schneider J, Arnrich B, Konigorski S. Comparison of {Bayesian Networks}, {G-estimation} and linear models to estimate causal treatment effects in aggregated {N-of-1} trials with carry-over effects. \emph{BMC Medical Research Methodology}. 2023;23(1):191. doi:\href{https://doi.org/10.1186/s12874-023-02012-5}{10.1186/s12874-023-02012-5}}

\bibitem[\citeproctext]{ref-raskind2003}
\CSLLeftMargin{18. }%
\CSLRightInline{{Raskind MA, Peskind ER, Kanter ED, et al.} Reduction of nightmares and other {PTSD} symptoms in combat veterans by prazosin: A placebo-controlled study. \emph{American Journal of Psychiatry}. 2003;160(2):371-373.}

\bibitem[\citeproctext]{ref-raskind2018}
\CSLLeftMargin{19. }%
\CSLRightInline{Raskind MA, Peskind ER, Chow B, et al. Trial of {Prazosin} for {Post-Traumatic Stress Disorder} in {Military Veterans}. \emph{New England Journal of Medicine}. 2018;378(6):507-517. doi:\href{https://doi.org/10.1056/NEJMoa1507598}{10.1056/NEJMoa1507598}}

\bibitem[\citeproctext]{ref-Khazaie2016Prazosin}
\CSLLeftMargin{20. }%
\CSLRightInline{Khazaie H, Nasouri M, Ghadami MR. Prazosin for {Trauma Nightmares} and {Sleep Disturbances} in {Combat Veterans} with {Post-Traumatic Stress Disorder}. \emph{Iranian Journal of Psychiatry and Behavioral Sciences}. 2016;10(3). doi:\href{https://doi.org/10.17795/ijpbs-2603}{10.17795/ijpbs-2603}}

\bibitem[\citeproctext]{ref-mayerRandomizedControlledPilot2023}
\CSLLeftMargin{21. }%
\CSLRightInline{Mayer CL, Savage PJ, Engle CK, et al. Randomized controlled pilot trial of prazosin for prophylaxis of posttraumatic headaches in active-duty service members and veterans. \emph{Headache: The Journal of Head and Face Pain}. 2023;63(6):751-762. doi:\href{https://doi.org/10.1111/head.14529}{10.1111/head.14529}}

\bibitem[\citeproctext]{ref-morris2019simulation}
\CSLLeftMargin{22. }%
\CSLRightInline{Morris TP, White IR, Crowther MJ. Using simulation studies to evaluate statistical methods. \emph{Statistics in Medicine}. 2019;38(11):2074-2102. doi:\href{https://doi.org/10.1002/sim.8086}{10.1002/sim.8086}}

\bibitem[\citeproctext]{ref-raskind2013}
\CSLLeftMargin{23. }%
\CSLRightInline{Raskind MA, Peterson K, Williams T, et al. A {Trial} of {Prazosin} for {Combat Trauma PTSD With Nightmares} in {Active-Duty Soldiers Returned From Iraq} and {Afghanistan}. \emph{American Journal of Psychiatry}. 2013;170(9):1003-1010. doi:\href{https://doi.org/10.1176/appi.ajp.2013.12081133}{10.1176/appi.ajp.2013.12081133}}

\bibitem[\citeproctext]{ref-schmidBayesianModelsNof12021}
\CSLLeftMargin{24. }%
\CSLRightInline{Schmid CH, Staudenmayer J. Bayesian models for {N-of-1} trials. \emph{Harvard Data Science Review}. 2021;3. doi:\href{https://doi.org/10.1162/99608f92.c06f5e3d}{10.1162/99608f92.c06f5e3d}}

\bibitem[\citeproctext]{ref-jiangDesignAnalysis2025}
\CSLLeftMargin{25. }%
\CSLRightInline{Jiang S, Arnold SE, Betensky RA. Design and analysis of n-of-1 trials that incorporate sequential monitoring. \emph{Statistics in Medicine}. 2025;44(15-17):e70177. doi:\href{https://doi.org/10.1002/sim.70177}{10.1002/sim.70177}}

\bibitem[\citeproctext]{ref-selukarFrameworkSequential2025}
\CSLLeftMargin{26. }%
\CSLRightInline{Selukar S, Prince DK, May S. A framework for sequential monitoring of individual n-of-1 trials and combining results across a series of sequentially monitored n-of-1 trials. \emph{Clinical Trials}. 2025;22(3):257-266. doi:\href{https://doi.org/10.1177/17407745241304284}{10.1177/17407745241304284}}

\bibitem[\citeproctext]{ref-mcdonaldMakingSwitchCase2019}
\CSLLeftMargin{27. }%
\CSLRightInline{McDonald S, Milne S, Beyene J, Sarti A, Thabane L. Making the switch: {From} case studies to {N-of-1} trials. \emph{Psychiatry Research}. 2019;278:86-93. doi:\href{https://doi.org/10.1016/j.psychres.2019.05.047}{10.1016/j.psychres.2019.05.047}}

\end{CSLReferences}

\protect\phantomsection\label{refs}
\begin{CSLReferences}{0}{1}
\bibitem[\citeproctext]{ref-lillie2011}
\CSLLeftMargin{1. }%
\CSLRightInline{Lillie EO, Patay B, Diamant J, Issell B, Topol EJ, Schork NJ. The n-of-1 clinical trial: The ultimate strategy for individualizing medicine? \emph{Personalized Medicine}. 2011;8(2):161-173. doi:\href{https://doi.org/10.2217/pme.11.7}{10.2217/pme.11.7}}

\bibitem[\citeproctext]{ref-nam2021expanding}
\CSLLeftMargin{2. }%
\CSLRightInline{National Academy of Medicine. Expanding the role of {N-of-1} trials in the precision medicine era: {Action} priorities and practical considerations. Published online 2021.}

\bibitem[\citeproctext]{ref-kravitzSinglepatientNof1Trials2014}
\CSLLeftMargin{3. }%
\CSLRightInline{{Kravitz RL, Duan N, Braslow J, et al.} Single-patient (n-of-1) trials: A pragmatic clinical decision methodology for patient-centered comparative effectiveness research. \emph{Journal of Clinical Epidemiology}. 2014;67(8):833-842. doi:\href{https://doi.org/10.1016/j.jclinepi.2013.04.006}{10.1016/j.jclinepi.2013.04.006}}

\bibitem[\citeproctext]{ref-marshallScopingReviewRandomized2022}
\CSLLeftMargin{4. }%
\CSLRightInline{Marshall S, Haywood K, Fitzpatrick R. A scoping review of randomized trials assessing the impact of n-of-1 trials on clinical outcomes. \emph{Journal of Clinical Medicine}. 2022;11(11):3026. doi:\href{https://doi.org/10.3390/jcm11113026}{10.3390/jcm11113026}}

\bibitem[\citeproctext]{ref-zucker2010}
\CSLLeftMargin{5. }%
\CSLRightInline{Zucker DR, Ruthazer R, Schmid CH. Individual ({N-of-1}) trials can be combined to give population comparative treatment effect estimates: Methodologic considerations. \emph{Journal of Clinical Epidemiology}. 2010;63(12):1312-1323. doi:\href{https://doi.org/10.1016/j.jclinepi.2010.04.020}{10.1016/j.jclinepi.2010.04.020}}

\bibitem[\citeproctext]{ref-chenComparisonAggregatedNof12020}
\CSLLeftMargin{6. }%
\CSLRightInline{Chen Y, Li P, Yang K. Comparison of aggregated {N-of-1} trials with parallel and crossover randomized controlled trials using simulation studies. \emph{PLoS One}. 2020;15(1):e0216949. doi:\href{https://doi.org/10.1371/journal.pone.0216949}{10.1371/journal.pone.0216949}}

\bibitem[\citeproctext]{ref-vohraCONSORTExtensionReporting2015}
\CSLLeftMargin{7. }%
\CSLRightInline{{Vohra S, Shamseer L, Sampson M, et al.} {CONSORT} extension for reporting {N-of-1} trials ({CENT}) 2015 {Statement}. \emph{Journal of Clinical Epidemiology}. 2015;76:9-17. doi:\href{https://doi.org/10.1016/j.jclinepi.2015.05.004}{10.1016/j.jclinepi.2015.05.004}}

\bibitem[\citeproctext]{ref-sennSampleSizeConsiderations2019}
\CSLLeftMargin{8. }%
\CSLRightInline{Senn S. Sample size considerations for n-of-1 trials. \emph{Statistical Methods in Medical Research}. 2019;28(2):372-383. doi:\href{https://doi.org/10.1177/0962280217726801}{10.1177/0962280217726801}}

\bibitem[\citeproctext]{ref-arendsSampleSizeConsiderations2019}
\CSLLeftMargin{9. }%
\CSLRightInline{{Arends RM, Hoekstra-Weebers JE, Sleijfer S, et al.} Sample size considerations for n-of-1 trials. \emph{Statistical Methods in Medical Research}. 2019;28(2):372-384. doi:\href{https://doi.org/10.1177/0962280217720357}{10.1177/0962280217720357}}

\bibitem[\citeproctext]{ref-yapBehavioralCarryoverEffect2024}
\CSLLeftMargin{10. }%
\CSLRightInline{Shi D, Ye T. Behavioral carry-over effect and power consideration in crossover trials. \emph{Biometrics}. 2024;80(2):ujae023. doi:\href{https://doi.org/10.1093/biomtc/ujae023}{10.1093/biomtc/ujae023}}

\bibitem[\citeproctext]{ref-tangTtestsNof1Trials2020}
\CSLLeftMargin{11. }%
\CSLRightInline{Tang W, Landes RD. Some t-tests for {N-of-1} trials with serial correlation. \emph{PLoS One}. 2020;15(2):e0228077. doi:\href{https://doi.org/10.1371/journal.pone.0228077}{10.1371/journal.pone.0228077}}

\bibitem[\citeproctext]{ref-wang2019}
\CSLLeftMargin{12. }%
\CSLRightInline{Wang Y, Schork NJ. Power and {Design Issues} in {Crossover-Based N-Of-1 Clinical Trials} with {Fixed Data Collection Periods}. \emph{Healthcare}. 2019;7(3):84. doi:\href{https://doi.org/10.3390/healthcare7030084}{10.3390/healthcare7030084}}

\bibitem[\citeproctext]{ref-liuConsiderationsCrossoverDesign2021}
\CSLLeftMargin{13. }%
\CSLRightInline{Liu J, Kong M, Zhou J. Considerations for crossover design in clinical study. \emph{Contemporary Clinical Trials Communications}. 2021;23:100818. doi:\href{https://doi.org/10.1016/j.conctc.2021.100818}{10.1016/j.conctc.2021.100818}}

\bibitem[\citeproctext]{ref-blackstonComparisonAggregatedNof12019}
\CSLLeftMargin{14. }%
\CSLRightInline{Blackston JW, Chapple AG, McGree JM, McDonald S, Nikles J. Comparison of {Aggregated N-of-1 Trials} with {Parallel} and {Crossover Randomized Controlled Trials Using Simulation Studies}. \emph{Healthcare}. 2019;7(4):137. doi:\href{https://doi.org/10.3390/healthcare7040137}{10.3390/healthcare7040137}}

\bibitem[\citeproctext]{ref-carrozzoTwoArm2025}
\CSLLeftMargin{15. }%
\CSLRightInline{Carrozzo AE, Zimmermann G, Bathke AC, Neunhaeuserer D, Niebauer J, Kulnik ST. Two-arm crossover randomized controlled trial versus meta-analysis of n-of-1 studies: Comparison of statistical efficiency in determining an intervention effect. \emph{Biometrical Journal}. 2025;67(2):e70045. doi:\href{https://doi.org/10.1002/bimj.70045}{10.1002/bimj.70045}}

\bibitem[\citeproctext]{ref-hendricksonOptimizing2020}
\CSLLeftMargin{16. }%
\CSLRightInline{Hendrickson RC, Thomas RG, Schork NJ, Raskind MA. Optimizing {Aggregated N-Of-1 Trial Designs} for {Predictive Biomarker Validation}: {Statistical Methods} and {Theoretical Findings}. \emph{Frontiers in Digital Health}. 2020;2. doi:\href{https://doi.org/10.3389/FDGTH.2020.00013/FULL}{10.3389/FDGTH.2020.00013/FULL}}

\bibitem[\citeproctext]{ref-gaertnerCarryover2023}
\CSLLeftMargin{17. }%
\CSLRightInline{Gärtner T, Schneider J, Arnrich B, Konigorski S. Comparison of {Bayesian Networks}, {G-estimation} and linear models to estimate causal treatment effects in aggregated {N-of-1} trials with carry-over effects. \emph{BMC Medical Research Methodology}. 2023;23(1):191. doi:\href{https://doi.org/10.1186/s12874-023-02012-5}{10.1186/s12874-023-02012-5}}

\bibitem[\citeproctext]{ref-raskind2003}
\CSLLeftMargin{18. }%
\CSLRightInline{{Raskind MA, Peskind ER, Kanter ED, et al.} Reduction of nightmares and other {PTSD} symptoms in combat veterans by prazosin: A placebo-controlled study. \emph{American Journal of Psychiatry}. 2003;160(2):371-373.}

\bibitem[\citeproctext]{ref-raskind2018}
\CSLLeftMargin{19. }%
\CSLRightInline{Raskind MA, Peskind ER, Chow B, et al. Trial of {Prazosin} for {Post-Traumatic Stress Disorder} in {Military Veterans}. \emph{New England Journal of Medicine}. 2018;378(6):507-517. doi:\href{https://doi.org/10.1056/NEJMoa1507598}{10.1056/NEJMoa1507598}}

\bibitem[\citeproctext]{ref-Khazaie2016Prazosin}
\CSLLeftMargin{20. }%
\CSLRightInline{Khazaie H, Nasouri M, Ghadami MR. Prazosin for {Trauma Nightmares} and {Sleep Disturbances} in {Combat Veterans} with {Post-Traumatic Stress Disorder}. \emph{Iranian Journal of Psychiatry and Behavioral Sciences}. 2016;10(3). doi:\href{https://doi.org/10.17795/ijpbs-2603}{10.17795/ijpbs-2603}}

\bibitem[\citeproctext]{ref-mayerRandomizedControlledPilot2023}
\CSLLeftMargin{21. }%
\CSLRightInline{Mayer CL, Savage PJ, Engle CK, et al. Randomized controlled pilot trial of prazosin for prophylaxis of posttraumatic headaches in active-duty service members and veterans. \emph{Headache: The Journal of Head and Face Pain}. 2023;63(6):751-762. doi:\href{https://doi.org/10.1111/head.14529}{10.1111/head.14529}}

\bibitem[\citeproctext]{ref-morris2019simulation}
\CSLLeftMargin{22. }%
\CSLRightInline{Morris TP, White IR, Crowther MJ. Using simulation studies to evaluate statistical methods. \emph{Statistics in Medicine}. 2019;38(11):2074-2102. doi:\href{https://doi.org/10.1002/sim.8086}{10.1002/sim.8086}}

\bibitem[\citeproctext]{ref-raskind2013}
\CSLLeftMargin{23. }%
\CSLRightInline{Raskind MA, Peterson K, Williams T, et al. A {Trial} of {Prazosin} for {Combat Trauma PTSD With Nightmares} in {Active-Duty Soldiers Returned From Iraq} and {Afghanistan}. \emph{American Journal of Psychiatry}. 2013;170(9):1003-1010. doi:\href{https://doi.org/10.1176/appi.ajp.2013.12081133}{10.1176/appi.ajp.2013.12081133}}

\bibitem[\citeproctext]{ref-schmidBayesianModelsNof12021}
\CSLLeftMargin{24. }%
\CSLRightInline{Schmid CH, Staudenmayer J. Bayesian models for {N-of-1} trials. \emph{Harvard Data Science Review}. 2021;3. doi:\href{https://doi.org/10.1162/99608f92.c06f5e3d}{10.1162/99608f92.c06f5e3d}}

\bibitem[\citeproctext]{ref-jiangDesignAnalysis2025}
\CSLLeftMargin{25. }%
\CSLRightInline{Jiang S, Arnold SE, Betensky RA. Design and analysis of n-of-1 trials that incorporate sequential monitoring. \emph{Statistics in Medicine}. 2025;44(15-17):e70177. doi:\href{https://doi.org/10.1002/sim.70177}{10.1002/sim.70177}}

\bibitem[\citeproctext]{ref-selukarFrameworkSequential2025}
\CSLLeftMargin{26. }%
\CSLRightInline{Selukar S, Prince DK, May S. A framework for sequential monitoring of individual n-of-1 trials and combining results across a series of sequentially monitored n-of-1 trials. \emph{Clinical Trials}. 2025;22(3):257-266. doi:\href{https://doi.org/10.1177/17407745241304284}{10.1177/17407745241304284}}

\bibitem[\citeproctext]{ref-mcdonaldMakingSwitchCase2019}
\CSLLeftMargin{27. }%
\CSLRightInline{McDonald S, Milne S, Beyene J, Sarti A, Thabane L. Making the switch: {From} case studies to {N-of-1} trials. \emph{Psychiatry Research}. 2019;278:86-93. doi:\href{https://doi.org/10.1016/j.psychres.2019.05.047}{10.1016/j.psychres.2019.05.047}}

\end{CSLReferences}

\begin{center}\rule{0.5\linewidth}{0.5pt}\end{center}

\textbf{Competing Interests:} The authors declare no competing
interests.

\textbf{Data Availability:} All simulation code and results are
available at the project repository. No human subjects data were
collected.

\textbf{Funding:} {[}To be filled in based on actual funding sources{]}

\end{document}
